% Standalone problem document, generated from the adjacent README.md.
% Compile with XeLaTeX. Mathematical content is maintained in Markdown.
\documentclass[11pt,a4paper]{article}
\usepackage[fontsize=10.5pt]{scrextend}
\usepackage[margin=22mm,headheight=15pt,headsep=9mm,footskip=11mm]{geometry}
\usepackage{fontspec}
\setmainfont{texgyrepagella}[Extension=.otf,UprightFont=*-regular,BoldFont=*-bold,ItalicFont=*-italic,BoldItalicFont=*-bolditalic]
\setsansfont{texgyreheros}[Extension=.otf,UprightFont=*-regular,BoldFont=*-bold,ItalicFont=*-italic,BoldItalicFont=*-bolditalic]
\setmonofont{texgyrecursor}[Extension=.otf,UprightFont=*-regular,BoldFont=*-bold,ItalicFont=*-italic,BoldItalicFont=*-bolditalic,Scale=MatchLowercase]
\usepackage{amsmath,amssymb}
\usepackage{unicode-math}
\setmathfont{texgyrepagella-math.otf}
\usepackage{xcolor,microtype,enumitem,needspace,fancyhdr}
\usepackage{bookmark}
\definecolor{ink}{HTML}{17344B}
\definecolor{accent}{HTML}{197D80}
\definecolor{muted}{HTML}{596773}
\hypersetup{colorlinks=true,linkcolor=accent,urlcolor=accent,pdftitle={SP-11 - The
delta conjecture for minimum symmetric
rank},pdfauthor={Open Problems in Numerical Linear Algebra}}
\urlstyle{same}
\setlength{\parindent}{0pt}
\setlength{\parskip}{5pt plus 1pt minus 1pt}
\linespread{1.04}
\setlength{\emergencystretch}{3em}
\widowpenalty=10000
\clubpenalty=10000
\displaywidowpenalty=10000
\allowdisplaybreaks[1]
\setlist{nosep,leftmargin=1.5em}
\providecommand{\tightlist}{\setlength{\itemsep}{0pt}\setlength{\parskip}{0pt}}
\providecommand{\pandocbounded}[1]{#1}
\pagestyle{fancy}
\fancyhf{}
\fancyhead[L]{\sffamily\footnotesize\color{muted}OPEN PROBLEMS IN NUMERICAL LINEAR ALGEBRA}
\fancyhead[R]{\sffamily\bfseries\footnotesize\color{accent}SP-11}
\fancyfoot[L]{\parbox[b]{0.9\textwidth}{\sffamily\fontsize{8}{10}\selectfont\color{muted}Editorial ratings; a bounded search cannot certify that a problem remains unsolved.\\Literature check: 2026-09-10}}
\fancyfoot[R]{\sffamily\footnotesize\color{muted}\thepage}
\renewcommand{\headrulewidth}{0.3pt}
\renewcommand{\headrule}{\hbox to\headwidth{\color{accent}\leaders\hrule height \headrulewidth\hfill}}
\makeatletter
\renewcommand{\section}{\Needspace{5\baselineskip}\@startsection{section}{1}{0pt}{12pt plus 2pt minus 2pt}{4pt}{\sffamily\bfseries\large\color{ink}}}
\renewcommand{\subsection}{\Needspace{4\baselineskip}\@startsection{subsection}{2}{0pt}{9pt plus 1pt minus 1pt}{3pt}{\sffamily\bfseries\normalsize\color{ink}}}
\makeatother
\setcounter{secnumdepth}{0}
\begin{document}
{\sffamily\small\color{muted}Eigenvalues and inverse problems\par}
\vspace{3pt}
{\raggedright\hyphenpenalty=10000\exhyphenpenalty=10000\sffamily\bfseries\fontsize{21}{24}\selectfont\color{ink}The
delta conjecture for minimum symmetric rank\par}
\vspace{7pt}
{\sffamily\small\textbf{Difficulty:} extreme\quad\textbf{Importance:} interesting
to the community\par}
{\sffamily\small\bfseries\color{accent}Status: Partially resolved\par}
\vspace{4pt}
{\color{accent}\hrule height 0.8pt}
\vspace{6pt}
\textbf{Rating rationale:} Extreme reflects the longstanding all-graph
gap between a local degree constraint and an exact matrix-realization
guarantee. Community impact comes from controlling achievable eigenvalue
multiplicity across symmetric sparsity patterns.

\subsection{Problem statement}\label{problem-statement}

Let \(G\) be any finite simple undirected graph on \(n\ge1\) vertices,
and let \(\delta(G)\) be its minimum vertex degree. Let
\(\mathcal S(G)\) consist of the real symmetric \(n\times n\) matrices
whose off-diagonal nonzero entries occur exactly at edges of \(G\), with
unrestricted diagonal entries.

Must there exist \(A\in\mathcal S(G)\) such that \[
\dim\ker A\ge\delta(G)?
\] Equivalently, with
\(\operatorname{mr}(G)=\min_{A\in\mathcal S(G)}\operatorname{rank}A\),
is \[
\operatorname{mr}(G)\le n-\delta(G)
\] valid for every \(G\)?

\subsection{Relevance and ratings}\label{relevance-and-ratings}

This asks whether local sparsity information guarantees feasibility of a
symmetric matrix with a specified large nullspace. It belongs to
structured inverse eigenvalue and low-rank matrix construction. It does
not prescribe the numerical edge weights, and does not require positive
semidefiniteness.

\subsection{References}\label{references}

\begin{itemize}
\tightlist
\item
  F. Barioli, S. M. Fallat, H. Gupta, and Z. Li, \emph{The weak version
  of the graph complement conjecture and partial results for the delta
  conjecture}, Discrete Mathematics 349 (2026), 114861, §1, paragraph
  immediately before Conjecture 1.5
  (\href{https://arxiv.org/html/2505.24577v1}{primary manuscript};
  \href{https://doi.org/10.1016/j.disc.2025.114861}{journal}).
\item
  S. M. Fallat and L. Hogben, \emph{The minimum rank of symmetric
  matrices described by a graph: A survey}, Linear Algebra and its
  Applications 426 (2007), 558--582, minimum-rank definitions and
  inverse eigenvalue interpretation
  (\href{https://aimath.org/WWN/matrixspectrum/FallatHogbenMinRank07.pdf}{primary
  manuscript};
  \href{https://doi.org/10.1016/j.laa.2007.05.036}{journal}).
\end{itemize}

\subsection{Status check}\label{status-check}

The 2026 article treats the original conjecture as open and establishes
results under additional graph restrictions. Its numbered Conjecture 1.5
is a stronger positive-semidefinite/SAP statement; the present entry
retains the original real symmetric formulation stated directly before
it. Searches for ``minimum rank'', ``delta conjecture'', ``proof'', and
2025--2026 found no general resolution. The stronger variants are not
counted as additional problems.

\subsection{Audit update --- 2026-09-10}\label{audit-update-2026-09-10}

\textbf{Partially resolved:} Barioli--Fallat--Gupta--Li, Theorem
2.11(a), proves the stronger bound \(\nu(G)\ge\delta(G)\) for graphs of
girth at least \(11\) and minimum degree at least \(4\). Parts (b)--(e)
give further explicit girth and forbidden-subgraph regimes. These
PSD/SAP witnesses also belong to the unrestricted symmetric class
displayed here. The full primary manuscript, 2026 publication record and
later-resolution searches were checked; no proof for all graphs was
located.
\end{document}
